% sections/related.tex -- Section 2 (paper_v2_vn, translated in L4 from paper_v2's
% sections/related.tex; terms as in glossary_vn.tex). Every work cited was read for this project;
% each claim was checked against the paper's text.
\section{Các nghiên cứu liên quan}
\label{sec:related}

\paragraph{Gợi ý dựa trên meta-path.} MCRec gán trọng số cho ngữ cảnh dựa trên meta-path
của từng cặp người dùng--sản phẩm bằng cơ chế đồng chú ý (co-attention) \citep{hu2018leveraging}.
Giống HAN \citep{wang2019heterogeneous}, HGRec \citep{shi2020hgrec}, FHRec \citep{mu2021fhrec} và
AMERec \citep{yan2021attention} học một embedding cho mỗi meta-path và hợp nhất
các embedding bằng cơ chế chú ý. Trong khi đó, metapath2vec học embedding của nút từ các bước đi ngẫu
nhiên theo một meta-path \citep{dong2017metapath2vec}. Tất cả các phương pháp này đều coi
meta-path là cho sẵn.

\paragraph{Phát hiện meta-path tự động.} Các phương pháp dựa trên truy vấn tìm $k$ meta-path tốt nhất nối hai nút cho trước \citep{zhu2020topk}. Các phương pháp dựa trên tìm kiếm
học xem một mô hình nên dùng những meta-path nào, hoặc những cấu trúc phong phú hơn. GEMS tiến hóa
các tập cấu trúc meta bằng thuật toán di truyền \citep{han2020genetic} và RMS mở
rộng các tập meta-path bằng học tăng cường \citep{ning2022automatic}. Cả hai đều đánh giá ứng
viên bằng chính mô hình gợi ý. Cách làm này rất chậm vì phải huấn luyện một mô hình gợi ý cho mỗi ứng viên: theo
so sánh của DiffMG, việc tìm kiếm của GEMS mất \val{diffmg-gems-min} đến \val{diffmg-gems-max}~phút~GPU
cho mỗi bộ dữ liệu, trong khi ba lần chạy tìm kiếm của DiffMG chỉ mất nhiều nhất
\val{diffmg-own-max}~phút \citep[Bảng~5]{ding2021diffmg}. Các phương pháp tìm kiếm theo
kiểu DARTS \citep{liu2019darts} giảm chi phí này bằng cách chỉ huấn luyện một lần: mọi ứng viên được
huấn luyện cùng lúc trong một siêu mạng. DiffMG tìm một đồ thị
meta dẫn dắt một mạng nơ-ron đồ thị cho bài toán phân loại nút và dự đoán liên kết
người dùng--sản phẩm, rồi giữ kết quả tốt nhất trong ba lần chạy tìm kiếm theo điểm xác thực
\citep{ding2021diffmg}. LMSPS lấy mẫu dần các meta-path từ một tập ứng viên lớn
\citep{li2024longrange}. Chúng tôi điều chỉnh LMSPS cho bài toán dự đoán liên kết. Khác với GEMS
và RMS, LMSPS đánh giá các tập của nó bằng mạng riêng, một mạng đóng vai trò thay cho HAN
(Mục~\ref{sec:res-weights} và~\ref{sec:conclusion}).

\paragraph{Đánh giá phương pháp tìm kiếm.} LMSPS so sánh các meta-path của nó với các meta-path
gốc của HAN, vốn được chọn thủ công, ngay trong HAN \citep[Bảng~4]{li2024longrange}, nhưng chỉ so
sánh với các meta-path ngẫu nhiên trong mạng riêng của nó, nơi các meta-path ngẫu nhiên đạt điểm thấp
hơn \citep[Bảng~1]{li2024longrange}. Một đối chứng ngẫu nhiên là cần thiết vì trong tìm kiếm kiến
trúc mạng nơ-ron, tìm kiếm ngẫu nhiên kết hợp dừng sớm đạt kết quả ít nhất ngang bằng ENAS, một
phương pháp tìm kiếm hàng đầu, trên hai bộ đánh giá chuẩn \citep{li2019random}. Xét trung bình, các
chiến lược tìm kiếm tiên tiến nhất cho kết quả tương tự chiến lược ngẫu nhiên
\citep{yu2020evaluating}. Trong số các phương pháp tìm kiếm meta-path, RMS vượt một đối chứng ngẫu
nhiên: trong cùng khoảng thời gian, đối chứng này huấn luyện mô hình gợi ý trên các tập chọn ngẫu
nhiên rồi giữ lại tập có điểm xác thực tốt nhất \citep{ning2022automatic}. GEMS và DiffMG không so
sánh với các cấu trúc ngẫu nhiên \citep{han2020genetic,ding2021diffmg}. Chúng tôi so sánh các tập
tìm được với các tập thủ công trên HAN và metapath2vec, nơi các meta-path được sử dụng, với các tập
chọn ngẫu nhiên từ chính các ứng viên của tìm kiếm làm đối chứng.

\paragraph{Đánh giá mô hình gợi ý.} Các mô hình cơ sở đơn giản thường vẫn cho kết quả cạnh tranh:
trong \val{dacrema-algos} mô hình gợi ý nơ-ron gần đây, chỉ kết quả của \val{dacrema-repro} mô hình
là tái hiện được với công sức hợp lý, trong đó \val{dacrema-beaten} mô hình thường có thể bị các
phương pháp heuristic đơn giản vượt qua \citep{dacrema2019progress}. Vì vậy, chúng tôi so sánh thêm
với mô hình phổ biến và một GAT không phân biệt kiểu, vốn là một mô hình cơ sở mạnh trên đồ thị
không đồng nhất \citep{lv2021hgb}. Thứ hạng của các mô hình theo độ đo lấy mẫu có thể khác với thứ
hạng khi xếp hạng toàn bộ \citep{krichene2022sampled}, nên chúng tôi cũng báo cáo kết quả xếp hạng
toàn bộ cho mọi tập meta-path.
